%%%%%%%%%%%%%%%%%%%%%%%%%%%%%%%%%%%%%%%%%%%%%%%%%%%%%%%%%%%%%%%%%%%%%%%%%%%%%
%% Documento base para escrita de monografias e relatórios do LabEPI.      %%
%% Aluno: Nome do Aluno <aluno@email.com>                                  %%
%% Orientador: Nome do Orientador <orientador@email.com>                   %%
%%%%%%%%%%%%%%%%%%%%%%%%%%%%%%%%%%%%%%%%%%%%%%%%%%%%%%%%%%%%%%%%%%%%%%%%%%%%%

%%%%%%%%%%%%%%%%%%%%%%%%%%%%%%%%%%%%%%%%%%%%%%%%%%%%%%%%%%%%%%%%%%%%%%%%%%%%%
%% DECLARAÇÂO DAS OPÇÔES E DO TIPO DO DOCUMENTO                            %%
%%%%%%%%%%%%%%%%%%%%%%%%%%%%%%%%%%%%%%%%%%%%%%%%%%%%%%%%%%%%%%%%%%%%%%%%%%%%%
% Definição dos possíveis tipos de documento
\def\doctypem{m}
\def\doctyper{r}

% Escolha do tipo de documento
\def\doctype{m} % Para monografias
%\def\doctype{r} % Para relatórios técnicos

\if\doctype\doctypem          % Se é do tipo monografia
\documentclass[twoside,a4paper,11pt]{book}
\fi
\if\doctype\doctyper           % Se é do tipo relatório
\documentclass[a4paper,11pt]{report}
\fi

%%%%%%%%%%%%%%%%%%%%%%%%%%%%%%%%%%%%%%%%%%%%%%%%%%%%%%%%%%%%%%%%%%%%%%%%%%%%%
%% PACOTES UTILIZADOS                                                      %%
%%%%%%%%%%%%%%%%%%%%%%%%%%%%%%%%%%%%%%%%%%%%%%%%%%%%%%%%%%%%%%%%%%%%%%%%%%%%%
%\usepackage[usenames,dvips]{color}
%\usepackage{ae}
%\usepackage{afterpage}
%\usepackage{appendix}
%\usepackage{eukdate}
%\usepackage{listings}
%\usepackage{lscape}
%\usepackage{multirow}
%\usepackage{nomencl}
%\usepackage{pslatex}
%\usepackage{rotating}
%\usepackage{setspace}
%\usepackage{supertabular}
%\usepackage{tabularx}
%\usepackage{tikz}
%\usepackage{url}
%\usepackage{verbatim}
%\usepackage{showframe}
\usepackage[top=3.5cm,bottom=3.5cm,left=3cm,right=3cm]{geometry}
%\usepackage[square]{natbib}
\usepackage{natbib}
\usepackage[labelfont=bf,font=small]{caption}
\usepackage[plain,chapter]{algorithm}
\usepackage[portuges,brazil]{babel}
\usepackage[utf8]{inputenc}
\usepackage{algorithmic}
\usepackage{amsmath}
\usepackage{amssymb}
\usepackage{amsthm}
\usepackage{booktabs}
\usepackage{paralist}
\usepackage{enumitem}
\usepackage{fancyhdr}
\usepackage{graphicx}
\usepackage{indentfirst}
\usepackage{lineno}
\usepackage{longtable}
\usepackage{thmtools}
\usepackage{lastpage}
\usepackage{imakeidx}
\usepackage[pagebackref]{hyperref}
\usepackage{titlesec}
\usepackage{pgfgantt}
\usepackage{fancyvrb}
\usepackage{emptypage}

%%%%%%%%%%%%%%%%%%%%%%%%%%%%%%%%%%%%%%%%%%%%%%%%%%%%%%%%%%%%%%%%%%%%%%%%%%%%%
%% PARÂMETROS                                                              %%
%%%%%%%%%%%%%%%%%%%%%%%%%%%%%%%%%%%%%%%%%%%%%%%%%%%%%%%%%%%%%%%%%%%%%%%%%%%%%

% Infenização
\hyphenation{net-works lay-out}

% Configuração do ambiente Verbatim do fancyvrb
\fvset{fontsize=\small,xleftmargin=1mm,xrightmargin=1mm}

% Configuração de diagramas de Gantt
\ganttset{%
  canvas/.append style={fill=none, dotted},
  %vgrid={draw={none}, dotted},
  vgrid,
  hgrid,
  today label=Prazo,
  group height=.3,
  group left shift=0,
  group right shift=0,
  group peaks height=1,
  group/.style={draw=black,fill=black},    
  bar height=.3,
  bar/.style={fill=black},
  progress label text={#1\pct},
  title height=1}

\DeclareMathOperator*{\argmin}{arg\,min}
\DeclareMathOperator*{\argmax}{arg\,max}

% Teoremas
\declaretheoremstyle[
  spaceabove=6pt,
  spacebelow=6pt,
  headfont=\normalfont\bfseries,
  notefont=\mdseries,
  notebraces={(}{)},
  bodyfont=\normalfont,
  postheadspace=1em,
  qed=$\square$,
  numberwithin=chapter
]{thmsty}

\usepackage{listings}
\usepackage{xcolor}

\newcounter{nalg}[chapter] % defines algorithm counter for chapter-level
\renewcommand{\thenalg}{\thechapter .\arabic{nalg}} %defines appearance of the algorithm counter
\DeclareCaptionLabelFormat{algocaption}{Algorithm \thenalg} % defines a new caption label as Algorithm x.y

\lstnewenvironment{pseudocode}[1][] %defines the algorithm listing environment
{   
    \refstepcounter{nalg} %increments algorithm number
    \captionsetup{labelformat=algocaption,labelsep=colon} %defines the caption setup for: it ises label format as the declared caption label above and makes label and caption text to be separated by a ':'
    \lstset{ %this is the stype
        mathescape=true,
        frame=tB,
        numbers=left, 
        numberstyle=\tiny,
        basicstyle=\scriptsize, 
        keywordstyle=\color{black}\bfseries\em,
        keywords={,input, output, return, datatype, function, in, if, else, foreach, while, begin, end, } %add the keywords you want, or load a language as Rubens explains in his comment above.
        numbers=left,
        xleftmargin=.04\textwidth,
        #1 % this is to add specific settings to an usage of this environment (for instnce, the caption and referable label)
    }
}
{}

\definecolor{mGreen}{rgb}{0,0.6,0}
\definecolor{mGray}{rgb}{0.5,0.5,0.5}
\definecolor{mPurple}{rgb}{0.58,0,0.82}
\definecolor{backgroundColour}{rgb}{0.95,0.95,0.92}

\lstdefinestyle{CStyle}{
    backgroundcolor=\color{backgroundColour},   
    commentstyle=\color{mGreen},
    keywordstyle=\color{magenta},
    numberstyle=\tiny\color{mGray},
    stringstyle=\color{mPurple},
    basicstyle=\footnotesize,
    breakatwhitespace=false,         
    breaklines=true,                 
    captionpos=b,                    
    keepspaces=true,                 
    numbers=left,                    
    numbersep=5pt,                  
    showspaces=false,                
    showstringspaces=false,
    showtabs=false,                  
    tabsize=2,
    language=C
}


\declaretheorem[style=thmsty,name=Algoritmo]{algoritmo}

\declaretheorem[style=thmsty,name=Definição]{definicao}
\declaretheorem[style=thmsty,name=Premissa]{premissa}

\declaretheorem[style=thmsty,name=Afirmação]{afirmacao}
\declaretheorem[style=thmsty,name=Observação]{observacao}
\declaretheorem[style=thmsty,name=Corolário]{corolario}
\declaretheorem[style=thmsty,name=Lema]{lema}
\declaretheorem[style=thmsty,name=Teorema]{teorema}

\declaretheorem[style=thmsty,name=Nota]{nota}

% Índice
\makeindex[intoc,options={-s labepi.ist}]

\hypersetup{colorlinks=true,linkcolor=blue,citecolor=blue}

\index{recorrencia@recorrência|see{recursividade}}
\index{recursividade|see{recorrência}}

%%%%%%%%%%%%%%%%%%%%%%%%%%%%%%%%%%%%%%%%%%%%%%%%%%%%%%%%%%%%%%%%%%%%%%%%%%%%%
%% INÍCIO DO DOCUMENTO                                                     %%
%%%%%%%%%%%%%%%%%%%%%%%%%%%%%%%%%%%%%%%%%%%%%%%%%%%%%%%%%%%%%%%%%%%%%%%%%%%%%

\begin{document}

%%%%%%%%%%%%%%%%%%%%%%%%%%%%%%%%%%%%%%%%%%%%%%%%%%%%%%%%%%%%%%%%%%%%%%%%%%%%%
%% INÍCIO DO ARQUIVO DE COMANDOS                                           %%
%%%%%%%%%%%%%%%%%%%%%%%%%%%%%%%%%%%%%%%%%%%%%%%%%%%%%%%%%%%%%%%%%%%%%%%%%%%%%

% Muda o estilo do cabeçalho dos capítulos so for um relatório
\if\doctype\doctyper           % Se é do tipo relatório
\titleformat{\chapter}{\normalfont\huge\bf}{\thechapter.}{20pt}{\huge\bf}
\fi

\newcommand{\limparpagina}{%
\if\doctype\doctypem           % Se é do tipo monografia
  \cleardoublepage
\fi
\if\doctype\doctyper           % Se é do tipo relatório
  \clearpage
\fi}


% Comandos próprios
\newcommand*\pct{\scalebox{.9}{\%}}

\newcommand{\citacao}[2]{%
\begin{flushright}
\footnotesize
\parbox{0.95\textwidth}{\flushright{\it``#1''\\ \scriptsize #2}}
\end{flushright}}

\newcommand{\definirtitulo}[1]{\newcommand{\titulo}{#1}}
\newcommand{\definirtitulocurto}[1]{\newcommand{\titulocurto}{#1}}

% Algoritmos
\floatname{algorithm}{Algoritmo}
\renewcommand{\listalgorithmname}{Lista de Algoritmos}

\renewcommand{\algorithmicindent}{2.0em}
\renewcommand{\algorithmicrequire}{\textbf{algoritmo}}
\renewcommand{\algorithmicensure}{\textbf{fim algoritmo}}
\renewcommand{\algorithmicreturn}{\textbf{retorne}}
\renewcommand{\algorithmicor}{\textbf{ou}}
\renewcommand{\algorithmicand}{\textbf{e}}
\renewcommand{\algorithmicend}{\textbf{fim}}
\renewcommand{\algorithmicif}{\textbf{se}}
\renewcommand{\algorithmicthen}{\textbf{ent\~ao}}
\renewcommand{\algorithmicfor}{\textbf{para}}
\renewcommand{\algorithmicforall}{\textbf{para cada}}
\renewcommand{\algorithmicrepeat}{\textbf{repita}}
\renewcommand{\algorithmicuntil}{\textbf{enquanto}}
\renewcommand{\algorithmicwhile}{\textbf{enquanto}}
\renewcommand{\algorithmicdo}{\textbf{fa\c{c}a}}

% Cabeçalho e rodapé
\renewcommand{\thefootnote}{[\roman{footnote}]}
\setlength{\headheight}{14.5pt}

\fancypagestyle{labepi}{%
\fancyhead{}
\fancyfoot{}
\if\doctype\doctypem           % Se é do tipo monografia
  \fancyhead[LO]{\bf\nouppercase \rightmark}
  \fancyhead[RO]{\bf\thepage}
  \fancyhead[LE]{\bf\thepage}
  \fancyhead[RE]{\bf\nouppercase \leftmark}
\fi
\if\doctype\doctyper           % Se é do tipo relatório
  \fancyhead[R]{\bf\thepage}
  \fancyhead[L]{\bf\titulocurto}
\fi
\renewcommand{\headrulewidth}{0pt}
\renewcommand{\footrulewidth}{0pt}}

\fancypagestyle{plain}{%
\fancyhead{}
\fancyfoot{}
\fancyhead[R]{\bf\thepage}
\fancyhead[L]{\bf\titulocurto}
\renewcommand{\headrulewidth}{0pt}
\renewcommand{\footrulewidth}{0pt}}

\renewcommand{\qedsymbol}{{\bf C.Q.D.}}

% Bibliografia
\newcommand{\backrefnotcited}{\relax} %(Not cited.)
\newcommand{\backrefcitedsingle}[1]{\\(Citado na página~#1)}
\newcommand{\backrefcitedmultiple}[1]{\\(Citado nas páginas~#1)}
\renewcommand{\backreftwosep}{ e~}
\renewcommand{\backreflastsep}{ e~}
\renewcommand*{\backref}[1]{}
\renewcommand*{\backrefalt}[4]{%
\ifcase #1
  \backrefnotcited
\or
  \backrefcitedsingle{#2}
\else
  \backrefcitedmultiple{#2}
\fi}

%%%%%%%%%%%%%%%%%%%%%%%%%%%%%%%%%%%%%%%%%%%%%%%%%%%%%%%%%%%%%%%%%%%%%%%%%%%%%
%% FIM DO ARQUIVO DE COMANDOS                                              %%
%%%%%%%%%%%%%%%%%%%%%%%%%%%%%%%%%%%%%%%%%%%%%%%%%%%%%%%%%%%%%%%%%%%%%%%%%%%%%
       % Comandos próprios

\definirtitulo{%
Relatório técnico sobre a disciplina de Estrutura de Dados}

\definirtitulocurto{Relatório sobre Estrutura de Dados}

% Numeração de linhas para auxiliar nas correções (comentar na versão final)
%\linenumbers

\pagestyle{empty}       % Páginas sem numeração e sem cabeçalhos

%%%%%%%%%%%%%%%%%%%%%%%%%%%%%%%%%%%%%%%%%%%%%%%%%%%%%%%%%%%%%%%%%%%%%%%%%%%%%
%% PÁGINA DE ROSTO                                                         %%
%%%%%%%%%%%%%%%%%%%%%%%%%%%%%%%%%%%%%%%%%%%%%%%%%%%%%%%%%%%%%%%%%%%%%%%%%%%%%
\begin{titlepage}

\begin{center}

\small

\parbox{0.33\textwidth}{%
\includegraphics[width=0.31\textwidth]{image/ufrn.eps}}%
\parbox{0.67\textwidth}{%
\begin{center}%
\textbf{%
Universidade Federal do Rio Grande do Norte -- UFRN\\
Centro de Ensino Superior do Seridó -- CERES\\
Departamento de Computação e Tecnologia -- DCT\\
Bacharelado em Sistemas de Informação -- BSI}
\end{center}}

\vfill

\LARGE

\textbf{\titulo}

\vfill

\Large

\textbf{Iago Gouveia Gurgel}

\vfill

\normalsize

Orientador: Prof. Dr. João Paulo Souza Medeiros

\vfill

\hfill
\if\doctype\doctypem           % Se é do tipo monografia
\parbox{0.5\linewidth}{%
{\bf Relatório Técnico} apresentado ao Curso de Bacharelado em Sistemas de
Informação como parte dos requisitos para aprovação na disciplina de
Estrutura de Dados.}
\fi
\if\doctype\doctyper           % Se é do tipo relatório
\parbox{0.5\linewidth}{%
{\bf Relatório Técnico} apresentado ao Curso de Bacharelado em Sistemas de
Informação como parte dos requisitos para aprovação na atividade de Estágio
Obrigatório.}
\fi

\vfill

\large

\includegraphics[width=0.23\textwidth]{image/labepi.eps} \\
Laboratório de Elementos do Processamento da Informação -- LabEPI \\
Caicó, RN, \today

\end{center}

\end{titlepage}
%%%%%%%%%%%%%%%%%%%%%%%%%%%%%%%%%%%%%%%%%%%%%%%%%%%%%%%%%%%%%%%%%%%%%%%%%%%%%
%% FIM DA FOLHA DE ROSTO                                                   %%
%%%%%%%%%%%%%%%%%%%%%%%%%%%%%%%%%%%%%%%%%%%%%%%%%%%%%%%%%%%%%%%%%%%%%%%%%%%%%
        % Folha de rosto
% \include{content/record}       % Ficha catalográfica
\if\doctype\doctypem           % Se é do tipo monografia
  %%%%%%%%%%%%%%%%%%%%%%%%%%%%%%%%%%%%%%%%%%%%%%%%%%%%%%%%%%%%%%%%%%%%%%%%%%%%%
%% FOLHA DE ASSINATURAS                                                    %%
%%%%%%%%%%%%%%%%%%%%%%%%%%%%%%%%%%%%%%%%%%%%%%%%%%%%%%%%%%%%%%%%%%%%%%%%%%%%%

%\begin{titlepage}
%
%\vfill
%
%\begin{center}
%\includegraphics[height=\textheight]{image/folha-de-rosto.eps}
%\end{center}
%
%\vfill
%
%\end{titlepage}

\begin{titlepage}

\begin{center}

\LARGE

\textbf{\titulo}

\vfill

\Large

\textbf{Iago Gouveia Gurgel}

\end{center}

\vfill

\noindent

Relatório aprovado em \today\ pelo examinador:

\begin{center}

\vspace{1.5cm}\rule{0.95\linewidth}{1pt}
\parbox{0.9\linewidth}{%
Prof. Dr. João Paulo Souza Medeiros (orientador) \dotfill\ DCT/UFRN}

\end{center}

\end{titlepage}
%%%%%%%%%%%%%%%%%%%%%%%%%%%%%%%%%%%%%%%%%%%%%%%%%%%%%%%%%%%%%%%%%%%%%%%%%%%%%
%% FIM DA FOLHA DE ASSINATURAS                                             %%
%%%%%%%%%%%%%%%%%%%%%%%%%%%%%%%%%%%%%%%%%%%%%%%%%%%%%%%%%%%%%%%%%%%%%%%%%%%%%
 % Folha de assinatura dos componentes da banca
  %%%%%%%%%%%%%%%%%%%%%%%%%%%%%%%%%%%%%%%%%%%%%%%%%%%%%%%%%%%%%%%%%%%%%%%%%%%%%
%% INÍCIO DA DEDICATÓRIA                                                   %%
%%%%%%%%%%%%%%%%%%%%%%%%%%%%%%%%%%%%%%%%%%%%%%%%%%%%%%%%%%%%%%%%%%%%%%%%%%%%%

\begin{titlepage}

\vspace*{\fill}

\citacao{%
There is no end to education.\\
It is not that you read a book, pass an examination, and finish with
education.\\
The whole of life, from the moment you are born to the moment you die, is a
process of learning.}
{Jiddu Krishnamurti}

\vspace*{\fill}

\end{titlepage}


%%%%%%%%%%%%%%%%%%%%%%%%%%%%%%%%%%%%%%%%%%%%%%%%%%%%%%%%%%%%%%%%%%%%%%%%%%%%%
%% FIM DA DEDICATÓRIA                                                      %%
%%%%%%%%%%%%%%%%%%%%%%%%%%%%%%%%%%%%%%%%%%%%%%%%%%%%%%%%%%%%%%%%%%%%%%%%%%%%%
 % Dedicatória
  % \include{content/thanks}     % Agradecimentos
\fi
%%%%%%%%%%%%%%%%%%%%%%%%%%%%%%%%%%%%%%%%%%%%%%%%%%%%%%%%%%%%%%%%%%%%%%%%%%%%%%
%% INÍCIO DOS RESUMOS                                                      %%
%%%%%%%%%%%%%%%%%%%%%%%%%%%%%%%%%%%%%%%%%%%%%%%%%%%%%%%%%%%%%%%%%%%%%%%%%%%%%

\chapter*{Resumo}
\thispagestyle{empty}

Este trabalho apresenta uma visão geral da disciplina de Estrutura de Dados
ministrada pelo DSc. João Paulo Souza Medeiros. Neste relatório, foi discutido
principalmente as atividades requeridas pelo Professor da disciplina e também,
apêndices com informações e algoritmos extras da curiosidade do aluno.

\vspace{1.5ex}

{\bf Palavras-chave}: 
Algoritmos;
Estruturas de Dados;

% ------------------------------------------------------------------------- %

\chapter*{Abstract}
\thispagestyle{empty}

This work presents an overview of the Data Structure discipline
taught by DSc. João Paulo Souza Medeiros. In this report, it was discussed
mainly the activities required by the subject teacher and also,
appendices with extra information and algorithms from the student's curiosity.

\vspace{1.5ex}

{\bf Keywords}:
Algorithms;
Data Structures;

%%%%%%%%%%%%%%%%%%%%%%%%%%%%%%%%%%%%%%%%%%%%%%%%%%%%%%%%%%%%%%%%%%%%%%%%%%%%%
%% FIM DOS RESUMOS                                                         %%
%%%%%%%%%%%%%%%%%%%%%%%%%%%%%%%%%%%%%%%%%%%%%%%%%%%%%%%%%%%%%%%%%%%%%%%%%%%%%
     % Resumo do trabalho em português e em inglês

\if\doctype\doctypem           % Se é do tipo monografia
\frontmatter
\fi

\pagestyle{labepi}

\tableofcontents

% Essas listas são opicionais nos relatórios, use-as quando a quantidade e a
% conveniência forem o caso.
%\limparpagina
%\renewcommand{\listtheoremname}{Lista de Algoritmos}
%\listoftheorems[ignoreall,show={algoritmo}]
%\addcontentsline{toc}{chapter}{\listtheoremname}

%\limparpagina
%\renewcommand{\listtheoremname}{Lista de Definições}
%\listoftheorems[ignoreall,show={definicao,premissa}]
%\addcontentsline{toc}{chapter}{\listtheoremname}

%\limparpagina
%\listoffigures
%\addcontentsline{toc}{chapter}{Lista de Figuras}

%\limparpagina
%\listoftables
%\addcontentsline{toc}{chapter}{Lista de Tabelas}

%\limparpagina
%\renewcommand{\listtheoremname}{Lista de Teoremas}
%\listoftheorems[ignoreall,show={corolario,lema,teorema}]
%\addcontentsline{toc}{chapter}{\listtheoremname}
%\newpage

% \include{content/glossary}

\if\doctype\doctypem           % Se é do tipo monografia
\mainmatter
\fi

%%%%%%%%%%%%%%%%%%%%%%%%%%%%%%%%%%%%%%%%%%%%%%%%%%%%%%%%%%%%%%%%%%%%%%%%%%%%%
%% CAPÍTULOS                                                               %%
%%%%%%%%%%%%%%%%%%%%%%%%%%%%%%%%%%%%%%%%%%%%%%%%%%%%%%%%%%%%%%%%%%%%%%%%%%%%%
%%%%%%%%%%%%%%%%%%%%%%%%%%%%%%%%%%%%%%%%%%%%%%%%%%%%%%%%%%%%%%%%%%%%%%%%%%%%%
%% INÍCIO CAPÍTULO                                                         %%
%%%%%%%%%%%%%%%%%%%%%%%%%%%%%%%%%%%%%%%%%%%%%%%%%%%%%%%%%%%%%%%%%%%%%%%%%%%%%

\chapter{Grafos e N-P Completeness}
\label{cap:4:conclusoes}

\citacao{%
If we can really understand the problem,\\
the answer will come out of it,\\
because the answer is not separate from the problem.}
{Jiddu Krishnamurti}

%%%%%%%%%%%%%%%%%%%%%%%%%%%%%%%%%%%%%%%%%%%%%%%%%%%%%%%%%%%%%%%%%%%%%%%%%%%%%
%% FIM CAPÍTULO                                                            %%
%%%%%%%%%%%%%%%%%%%%%%%%%%%%%%%%%%%%%%%%%%%%%%%%%%%%%%%%%%%%%%%%%%%%%%%%%%%%%
 % Introdução (apresentação do trabalho)
%%%%%%%%%%%%%%%%%%%%%%%%%%%%%%%%%%%%%%%%%%%%%%%%%%%%%%%%%%%%%%%%%%%%%%%%%%%%%
%% INÍCIO CAPÍTULO                                                         %%
%%%%%%%%%%%%%%%%%%%%%%%%%%%%%%%%%%%%%%%%%%%%%%%%%%%%%%%%%%%%%%%%%%%%%%%%%%%%%

\chapter{Grafos e N-P Completeness}
\label{cap:4:conclusoes}

\citacao{%
If we can really understand the problem,\\
the answer will come out of it,\\
because the answer is not separate from the problem.}
{Jiddu Krishnamurti}

%%%%%%%%%%%%%%%%%%%%%%%%%%%%%%%%%%%%%%%%%%%%%%%%%%%%%%%%%%%%%%%%%%%%%%%%%%%%%
%% FIM CAPÍTULO                                                            %%
%%%%%%%%%%%%%%%%%%%%%%%%%%%%%%%%%%%%%%%%%%%%%%%%%%%%%%%%%%%%%%%%%%%%%%%%%%%%%
 % Análise de algoritmos de Ordenação
%%%%%%%%%%%%%%%%%%%%%%%%%%%%%%%%%%%%%%%%%%%%%%%%%%%%%%%%%%%%%%%%%%%%%%%%%%%%%
%% INÍCIO CAPÍTULO                                                         %%
%%%%%%%%%%%%%%%%%%%%%%%%%%%%%%%%%%%%%%%%%%%%%%%%%%%%%%%%%%%%%%%%%%%%%%%%%%%%%

\chapter{Grafos e N-P Completeness}
\label{cap:4:conclusoes}

\citacao{%
If we can really understand the problem,\\
the answer will come out of it,\\
because the answer is not separate from the problem.}
{Jiddu Krishnamurti}

%%%%%%%%%%%%%%%%%%%%%%%%%%%%%%%%%%%%%%%%%%%%%%%%%%%%%%%%%%%%%%%%%%%%%%%%%%%%%
%% FIM CAPÍTULO                                                            %%
%%%%%%%%%%%%%%%%%%%%%%%%%%%%%%%%%%%%%%%%%%%%%%%%%%%%%%%%%%%%%%%%%%%%%%%%%%%%%
 % Análise de algoritmos de Árvores e Tabelas de Dispersão
% %%%%%%%%%%%%%%%%%%%%%%%%%%%%%%%%%%%%%%%%%%%%%%%%%%%%%%%%%%%%%%%%%%%%%%%%%%%%%
%% INÍCIO CAPÍTULO                                                         %%
%%%%%%%%%%%%%%%%%%%%%%%%%%%%%%%%%%%%%%%%%%%%%%%%%%%%%%%%%%%%%%%%%%%%%%%%%%%%%

\chapter{Grafos e N-P Completeness}
\label{cap:4:conclusoes}

\citacao{%
If we can really understand the problem,\\
the answer will come out of it,\\
because the answer is not separate from the problem.}
{Jiddu Krishnamurti}

%%%%%%%%%%%%%%%%%%%%%%%%%%%%%%%%%%%%%%%%%%%%%%%%%%%%%%%%%%%%%%%%%%%%%%%%%%%%%
%% FIM CAPÍTULO                                                            %%
%%%%%%%%%%%%%%%%%%%%%%%%%%%%%%%%%%%%%%%%%%%%%%%%%%%%%%%%%%%%%%%%%%%%%%%%%%%%%
 % Grafos
% %%%%%%%%%%%%%%%%%%%%%%%%%%%%%%%%%%%%%%%%%%%%%%%%%%%%%%%%%%%%%%%%%%%%%%%%%%%%%
%% INÍCIO CAPÍTULO                                                         %%
%%%%%%%%%%%%%%%%%%%%%%%%%%%%%%%%%%%%%%%%%%%%%%%%%%%%%%%%%%%%%%%%%%%%%%%%%%%%%

\chapter{Grafos e N-P Completeness}
\label{cap:4:conclusoes}

\citacao{%
If we can really understand the problem,\\
the answer will come out of it,\\
because the answer is not separate from the problem.}
{Jiddu Krishnamurti}

%%%%%%%%%%%%%%%%%%%%%%%%%%%%%%%%%%%%%%%%%%%%%%%%%%%%%%%%%%%%%%%%%%%%%%%%%%%%%
%% FIM CAPÍTULO                                                            %%
%%%%%%%%%%%%%%%%%%%%%%%%%%%%%%%%%%%%%%%%%%%%%%%%%%%%%%%%%%%%%%%%%%%%%%%%%%%%%
 % Conclusão

%%%%%%%%%%%%%%%%%%%%%%%%%%%%%%%%%%%%%%%%%%%%%%%%%%%%%%%%%%%%%%%%%%%%%%%%%%%%%
%% APÊNDICE, BILIOGRAFIA E ÍNDICE REMISSIVO                                %%
%%%%%%%%%%%%%%%%%%%%%%%%%%%%%%%%%%%%%%%%%%%%%%%%%%%%%%%%%%%%%%%%%%%%%%%%%%%%%
\appendix
%%%%%%%%%%%%%%%%%%%%%%%%%%%%%%%%%%%%%%%%%%%%%%%%%%%%%%%%%%%%%%%%%%%%%%%%%%%%%%
%% INÍCIO CAPÍTULO                                                         %%
%%%%%%%%%%%%%%%%%%%%%%%%%%%%%%%%%%%%%%%%%%%%%%%%%%%%%%%%%%%%%%%%%%%%%%%%%%%%%

\chapter{Grafos e N-P Completeness}
\label{cap:4:conclusoes}

\citacao{%
If we can really understand the problem,\\
the answer will come out of it,\\
because the answer is not separate from the problem.}
{Jiddu Krishnamurti}

%%%%%%%%%%%%%%%%%%%%%%%%%%%%%%%%%%%%%%%%%%%%%%%%%%%%%%%%%%%%%%%%%%%%%%%%%%%%%
%% FIM CAPÍTULO                                                            %%
%%%%%%%%%%%%%%%%%%%%%%%%%%%%%%%%%%%%%%%%%%%%%%%%%%%%%%%%%%%%%%%%%%%%%%%%%%%%%

%%%%%%%%%%%%%%%%%%%%%%%%%%%%%%%%%%%%%%%%%%%%%%%%%%%%%%%%%%%%%%%%%%%%%%%%%%%%%%
%% INÍCIO CAPÍTULO                                                         %%
%%%%%%%%%%%%%%%%%%%%%%%%%%%%%%%%%%%%%%%%%%%%%%%%%%%%%%%%%%%%%%%%%%%%%%%%%%%%%

\chapter{Grafos e N-P Completeness}
\label{cap:4:conclusoes}

\citacao{%
If we can really understand the problem,\\
the answer will come out of it,\\
because the answer is not separate from the problem.}
{Jiddu Krishnamurti}

%%%%%%%%%%%%%%%%%%%%%%%%%%%%%%%%%%%%%%%%%%%%%%%%%%%%%%%%%%%%%%%%%%%%%%%%%%%%%
%% FIM CAPÍTULO                                                            %%
%%%%%%%%%%%%%%%%%%%%%%%%%%%%%%%%%%%%%%%%%%%%%%%%%%%%%%%%%%%%%%%%%%%%%%%%%%%%%

%%%%%%%%%%%%%%%%%%%%%%%%%%%%%%%%%%%%%%%%%%%%%%%%%%%%%%%%%%%%%%%%%%%%%%%%%%%%%%
%% INÍCIO CAPÍTULO                                                         %%
%%%%%%%%%%%%%%%%%%%%%%%%%%%%%%%%%%%%%%%%%%%%%%%%%%%%%%%%%%%%%%%%%%%%%%%%%%%%%

\chapter{Grafos e N-P Completeness}
\label{cap:4:conclusoes}

\citacao{%
If we can really understand the problem,\\
the answer will come out of it,\\
because the answer is not separate from the problem.}
{Jiddu Krishnamurti}

%%%%%%%%%%%%%%%%%%%%%%%%%%%%%%%%%%%%%%%%%%%%%%%%%%%%%%%%%%%%%%%%%%%%%%%%%%%%%
%% FIM CAPÍTULO                                                            %%
%%%%%%%%%%%%%%%%%%%%%%%%%%%%%%%%%%%%%%%%%%%%%%%%%%%%%%%%%%%%%%%%%%%%%%%%%%%%%

%%%%%%%%%%%%%%%%%%%%%%%%%%%%%%%%%%%%%%%%%%%%%%%%%%%%%%%%%%%%%%%%%%%%%%%%%%%%%%
%% INÍCIO CAPÍTULO                                                         %%
%%%%%%%%%%%%%%%%%%%%%%%%%%%%%%%%%%%%%%%%%%%%%%%%%%%%%%%%%%%%%%%%%%%%%%%%%%%%%

\chapter{Grafos e N-P Completeness}
\label{cap:4:conclusoes}

\citacao{%
If we can really understand the problem,\\
the answer will come out of it,\\
because the answer is not separate from the problem.}
{Jiddu Krishnamurti}

%%%%%%%%%%%%%%%%%%%%%%%%%%%%%%%%%%%%%%%%%%%%%%%%%%%%%%%%%%%%%%%%%%%%%%%%%%%%%
%% FIM CAPÍTULO                                                            %%
%%%%%%%%%%%%%%%%%%%%%%%%%%%%%%%%%%%%%%%%%%%%%%%%%%%%%%%%%%%%%%%%%%%%%%%%%%%%%

%%%%%%%%%%%%%%%%%%%%%%%%%%%%%%%%%%%%%%%%%%%%%%%%%%%%%%%%%%%%%%%%%%%%%%%%%%%%%%
%% INÍCIO CAPÍTULO                                                         %%
%%%%%%%%%%%%%%%%%%%%%%%%%%%%%%%%%%%%%%%%%%%%%%%%%%%%%%%%%%%%%%%%%%%%%%%%%%%%%

\chapter{Grafos e N-P Completeness}
\label{cap:4:conclusoes}

\citacao{%
If we can really understand the problem,\\
the answer will come out of it,\\
because the answer is not separate from the problem.}
{Jiddu Krishnamurti}

%%%%%%%%%%%%%%%%%%%%%%%%%%%%%%%%%%%%%%%%%%%%%%%%%%%%%%%%%%%%%%%%%%%%%%%%%%%%%
%% FIM CAPÍTULO                                                            %%
%%%%%%%%%%%%%%%%%%%%%%%%%%%%%%%%%%%%%%%%%%%%%%%%%%%%%%%%%%%%%%%%%%%%%%%%%%%%%

%%%%%%%%%%%%%%%%%%%%%%%%%%%%%%%%%%%%%%%%%%%%%%%%%%%%%%%%%%%%%%%%%%%%%%%%%%%%%%
%% INÍCIO CAPÍTULO                                                         %%
%%%%%%%%%%%%%%%%%%%%%%%%%%%%%%%%%%%%%%%%%%%%%%%%%%%%%%%%%%%%%%%%%%%%%%%%%%%%%

\chapter{Grafos e N-P Completeness}
\label{cap:4:conclusoes}

\citacao{%
If we can really understand the problem,\\
the answer will come out of it,\\
because the answer is not separate from the problem.}
{Jiddu Krishnamurti}

%%%%%%%%%%%%%%%%%%%%%%%%%%%%%%%%%%%%%%%%%%%%%%%%%%%%%%%%%%%%%%%%%%%%%%%%%%%%%
%% FIM CAPÍTULO                                                            %%
%%%%%%%%%%%%%%%%%%%%%%%%%%%%%%%%%%%%%%%%%%%%%%%%%%%%%%%%%%%%%%%%%%%%%%%%%%%%%


\if\doctype\doctypem           % Se é do tipo monografia
\backmatter
\fi

\limparpagina
%\nocite{*}
\bibliographystyle{labepi}
\bibliography{content/bibliography}
\addcontentsline{toc}{chapter}{\bibname}

\limparpagina
%\markboth{Índice Remissivo}{Índice Remissivo}
%\printindex

\end{document}

%%%%%%%%%%%%%%%%%%%%%%%%%%%%%%%%%%%%%%%%%%%%%%%%%%%%%%%%%%%%%%%%%%%%%%%%%%%%%
%% FIM DO DOCUMENTO                                                        %%
%%%%%%%%%%%%%%%%%%%%%%%%%%%%%%%%%%%%%%%%%%%%%%%%%%%%%%%%%%%%%%%%%%%%%%%%%%%%%
